\chapter{Themenregister über alle Bände}
\markboth{Themenregister}{Themenregister}

Dieses Register nennt Begriffe, die in mehreren Bänden vorkommen, und sagt, wo sie behandelt werden. Fett steht der Band, in dem der Begriff eingeführt wird; die übrigen Nennungen setzen ihn voraus oder greifen ihn wieder auf.

\begin{merksatz}
\boxtitel{Begriffe, die die Reihe durchziehen}
\begin{center}
{\small
\renewcommand{\arraystretch}{1.36}
\begin{longtable}{@{}P{4.2cm}P{4.6cm}@{}}
\toprule
\textbf{Begriff} & \textbf{Bände}\\
\midrule
\endhead
Abstraktionsschicht & \textbf{0}, 4, 5, 6, 21\\
Amdahlsches Gesetz & \textbf{4}, 5\\
Aufwandsklassen & \textbf{1}, 3, 11, 13\\
Automat, endlicher & \textbf{0}, 2, 6\\
Cache, Zwischenspeicher & \textbf{4}, 10, 23, 24\\
Entropie & \textbf{7}, 8, 13\\
Fehlerkorrektur & 7, 9, \textbf{15}\\
Fließbandverarbeitung & \textbf{4}, 5\\
Gleitkommazahl & \textbf{0}, 13, 21, 24\\
Halteproblem & \textbf{2}, 3, Begleitband\\
Idempotenz & 9, \textbf{23}, 24\\
Konsens & \textbf{9}, 10\\
Naht, Schnittstelle & 9, 17, \textbf{23}, 24\\
Nebenläufigkeit & \textbf{5}, 9, 22, 23\\
NP-Vollständigkeit & 2, \textbf{3}, 12\\
Reduktion & 2, \textbf{3}\\
Selbstbezug & 0, 2, \textbf{Begleitband}\\
Speicherhierarchie & \textbf{4}, 5, 10\\
Transaktion & 9, \textbf{10}\\
Turingmaschine & \textbf{2}, 3, Begleitband\\
Undefiniertes Verhalten & 6, \textbf{21}, 24\\
Universelle Maschine & \textbf{0}, 2, Begleitband\\
Vertrauensgrenze & 8, \textbf{19}, 23, 24\\
Verklemmung & \textbf{5}, 22\\
Wettlauf & \textbf{5}, 22, 24\\
Zeit, Uhren & \textbf{9}, 23, 24\\
Zustandsraum & 2, \textbf{12}, 15\\
\bottomrule
\end{longtable}}
\end{center}
\end{merksatz}

\begin{technik}
\boxtitel{Wo dieselbe Frage mehrfach gestellt wird}
\begin{center}
{\small
\renewcommand{\arraystretch}{1.45}
\begin{tabular}{@{}P{4.0cm}P{4.8cm}@{}}
\toprule
\textbf{Frage} & \textbf{Bände}\\
\midrule
Was kostet eine Abstraktion? & 0, 4, 6, 21, 23\\
Was kann ein System über sich selbst wissen? & 0, 2, 15, Begleitband\\
Wie unterscheidet man Rauschen von Signal? & 7, 13, 14\\
Wem darf man vertrauen? & 8, 9, 19, 23\\
Wann greift eine Regel nicht? & 17, 20, 22, \textbf{24}\\
\bottomrule
\end{tabular}}
\end{center}

\begin{center}
\begin{minipage}{0.86\textwidth}
\kursiv{Die letzte Zeile ist erst mit den Programmierbänden dazugekommen und ist die einzige Frage der Reihe, auf die es keine allgemeine Antwort gibt --- weshalb Band 24 aus zwanzig einzelnen Antworten besteht.}
\end{minipage}
\end{center}
\end{technik}

\begin{merksatz}
\boxtitel{Verweise zwischen den Bänden}
\begin{center}
{\small
\renewcommand{\arraystretch}{1.45}
\begin{tabular}{@{}P{2.4cm}P{6.4cm}@{}}
\toprule
\textbf{Band 21} & greift auf Band 0 (Zahlen), 4 (Speicher) und 6 (Übersetzer) zurück\\
\textbf{Band 22} & greift auf Band 5 (Nebenläufigkeit) zurück\\
\textbf{Band 23} & greift auf Band 9 (verteilte Systeme) und 19 (Sicherheit) zurück\\
\textbf{Band 24} & prüft Regeln aus den Bänden 17, 19, 20, 21, 22 und 23 am Einzelfall\\
\bottomrule
\end{tabular}}
\end{center}

\begin{center}
\begin{minipage}{0.86\textwidth}
\kursiv{Keiner dieser Verweise ist zwingend: Jeder Band erklärt, was er braucht. Die Verweise sagen nur, wo dasselbe ausführlicher steht.}
\end{minipage}
\end{center}
\end{merksatz}
